\section{The retained primary cyclic lattice at its specified coefficient prime}
\label{sec:tcl}

This calculation uses the ordered primary tensor and full Taylor unit
of TP.1 and TP.23--25. It computes its original cyclic map over the
specified coefficient prime, including the map before reduction, its
primitive containing lattice, its cokernel, and the exact reduction
sequence. The support carrier used below is the original
\(G_{C_n}\) of MFC.29 and MFC.35--44.

\subsection{The actual coefficient localization and the unchanged coordinates}

Use the single-primary coefficient specialization already specified
in SPF.1. Denote its retained coefficient subring of \(\mathbb C\)
by \(A\), its map by \(\alpha:A\to k_0\), and its characteristic
by the prime \(p\). Retain every coefficient of the full Taylor
class and its inverse, and put
\[
 \mathfrak q=\ker\alpha,\qquad S=A_{\mathfrak q},\qquad
 \widetilde\alpha:S\longrightarrow k_0,\qquad
 \widetilde\alpha(a/b)=\alpha(a)/\alpha(b).
 \tag{TCL.1}\label{eq:tcl-ring}
\]
Thus \(b\notin\mathfrak q\) in this formula. The image of
\(\alpha\) is a finite subring of a field. For each nonzero element
of that image, multiplication on the finite image is injective and
therefore surjective; its inverse lies in the image. The image is
consequently a field, proving that \(\mathfrak q\) is maximal.
The fraction formula is independent of representatives: equality
of two fractions gives equality after cross multiplication in
the domain \(A\), and applying \(\alpha\) and dividing by the
nonzero denominator images gives the asserted equality.
It preserves the ring operations by the same calculation.
The map extends the specified map, rather than selecting a new one.

Since \(A\subset\mathbb C\), both \(A\) and its localization
\(S\) are domains of characteristic zero. The element \(p\) is
a nonzero nonunit: its image is zero whereas the image of the unit
is one. Every integer prime to \(p\) is outside \(\mathfrak q\)
and is a unit in \(S\). In particular \(\mathbb Z_{(p)}\to S\)
is a well-defined injective scalar map. No assertion that \(S\)
is a discrete valuation ring, a principal ideal domain, or a ring
with all ideals in a chain is made or used. The algebra below uses
\(p>m\); the retained SPF specialization has the stronger
\(p>m+1\), as required there by its phase denominators.

Keep the selected primary coordinate \(\rho\), its full order
\(m\geq1\), and the ordered number of factors \(k\geq1\). Put
\[
 \begin{split}
 K&=k(m-1),\qquad y_i=s_i-\rho,\qquad
 E_S=S[s_1,\ldots,s_k]/((s_i-\rho)^m)_{i=1}^k
       =S[y_1,\ldots,y_k]/(y_i^m)_{i=1}^k,\\
 J&=\sum_{i=1}^k y_i,\qquad
 M=\sum_{i=1}^k s_i=k\rho+J,\qquad
 \mathcal C_S=S[T]/(T^{K+1}),\qquad
 \mathcal C_S^{\rm orig}=S[Z]/((Z-k\rho)^{K+1}).
 \end{split}
 \tag{TCL.2}\label{eq:tcl-objects}
\]
The two displayed presentations of \(E_S\) are identified by the
mutually inverse substitutions \(s_i=y_i+\rho\) and
\(y_i=s_i-\rho\). The cyclic-source coordinate map
\(\mathcal C_S^{\rm orig}\to\mathcal C_S\) is
\(Z\mapsto k\rho+T\), with inverse \(T\mapsto Z-k\rho\).
These substitutions respect the defining ideals and their composites
are the identity on each generator. They preserve the original objects
and will be recorded explicitly in the original monomial matrices.

Let
\(\mathcal B=\{0,\ldots,m-1\}^k\), ordered lexicographically.
Repeated division by the monic relations \(y_i^m\) proves that
\(\{y^{\mathbf b}:\mathbf b\in\mathcal B\}\) is an \(S\)-basis:
deleting monomials with any exponent at least \(m\) gives a
remainder, and the relation ideal is spanned by precisely the deleted
monomials, so no relation between retained monomials is possible.
Every surviving degree is at most \(K\); hence \(J^{K+1}=0\).

Write the retained one-factor unit and its retained inverse as
\[
 \begin{split}
 v(y)&=\sum_{j=0}^{m-1}a_jy^j,\qquad
 a_j=\frac{(g/h)^{(j)}(\rho)}{j!}
     =\frac{g^{(m+j)}(\rho)}{(m+j)!},\qquad h=(s-\rho)^m,\\
 w(y)&=\sum_{j=0}^{m-1}b_jy^j,\qquad
 b_0=a_0^{-1},\qquad
 b_r=-a_0^{-1}\sum_{j=1}^{r}a_jb_{r-j},\\
 U&=\prod_{i=1}^kv(y_i),\qquad
 U^{-1}=\prod_{i=1}^kw(y_i)\quad\hbox{in }E_S.
 \end{split}
 \tag{TCL.3}\label{eq:tcl-full-unit}
\]
These are the actual retained coefficients, all lying in \(A\)
by its specified construction. The derivative identity follows by
expanding \(g=(s-\rho)^m(g/h)\) at \(\rho\). The recurrence makes
the constant coefficient of \(v w\) equal to one and its coefficients
of degrees \(1,\ldots,m-1\) equal to zero. Thus \(v w=1\) modulo
\(y^m\), proving the inverse in every factor and then in the tensor.
In particular multiplication by \(U\) is an actual \(S\)-linear
automorphism of \(E_S\) commuting with multiplication by \(J\)
and \(M\). No coefficient of this unit has been discarded.

The unchanged unweighted cyclic ring map and the original unit-weighted
cyclic coefficient map are respectively
\[
 \begin{split}
 f:\mathcal C_S\longrightarrow E_S,&\qquad f(T^r)=J^r,\\
 f_U:\mathcal C_S\longrightarrow E_S,&\qquad f_U(T^r)=U J^r,\\
 f_U^{\rm orig}:\mathcal C_S^{\rm orig}\longrightarrow E_S,
 &\qquad f_U^{\rm orig}(Z^j)=U M^j.
 \end{split}
 \tag{TCL.4}\label{eq:tcl-original-map}
\]
The relation \(J^{K+1}=0\) proves that all maps are well-defined.
The last two are linear maps, and are maps of modules over the cyclic
source ring acting on the target by \(T\mapsto J\), respectively
\(Z\mapsto M\). They are not asserted to be unital ring maps:
their value at one is the retained \(U\). Their exact intertwining
identities are
\[
 f_U(Tc)=Jf_U(c),\qquad
 f_U^{\rm orig}(Zc)=M f_U^{\rm orig}(c).
 \tag{TCL.5}\label{eq:tcl-intertwining}
\]
Both follow by multiplication in the commutative algebra \(E_S\).
This is the coefficient realization of the original source map
\(\mathcal J_k\mathcal T_k\) in TP.25, restricted to polynomials
in the sum of its unchanged \(s_i\) coordinates.

\subsection{The primitive lattice and the exact factorial inclusion}

For \(0\leq r\leq K\), define
\[
 \mathcal B_r=\{\mathbf b\in\mathcal B:|\mathbf b|=r\},\qquad
 \mathbf b!=\prod_{i=1}^k b_i!,\qquad
 V_r=\sum_{\mathbf b\in\mathcal B_r}\frac{y^{\mathbf b}}{\mathbf b!},
 \qquad V_{K+1}=0.
 \tag{TCL.6}\label{eq:tcl-divided}
\]
Every denominator is a unit because \(b_i<m<p\). Each set
\(\mathcal B_r\) is nonempty. For \(m>1\) a specified member
\(\mathbf b^*(r)\) is obtained recursively by
\(b_i^*(r)=\min(m-1,r-\sum_{j<i}b_j^*(r))\).
The remaining degree is nonnegative at each step and is zero after
\(k\) steps because \(r\leq k(m-1)\). For \(m=1\) there is only
\(r=0\), with \(\mathbf b^*(0)=(0,\ldots,0)\).

The multinomial expansion and direct multiplication give
\[
 J^r=r!V_r\quad(0\leq r\leq K),\qquad
 J V_r=(r+1)V_{r+1}\quad(0\leq r\leq K).
 \tag{TCL.7}\label{eq:tcl-factorial}
\]
To prove the first formula, expand the product of \(r\) copies
of \(\sum_i y_i\). For a prescribed exponent vector \(\mathbf b\)
there are \(r!/\mathbf b!\) assignments of its \(r\) ordered
factors to the variables. Exponents at least \(m\) give zero in
the quotient. This leaves exactly the displayed coefficients.
For the second formula, the coefficient of a retained
\(y^{\mathbf b}\) of degree \(r+1\) is
\(\sum_{i:b_i>0}1/(b_i-1)!\prod_{j\ne i}1/b_j!
=(\sum_i b_i)/\mathbf b!=(r+1)/\mathbf b!\).
For \(r=K\) every resulting monomial is zero, which agrees with
the convention for \(V_{K+1}\).

Define coefficient functionals, a containing lattice and its projection by
\[
 \begin{split}
 \lambda_r(x)&=\mathbf b^*(r)!\,[y^{\mathbf b^*(r)}]x,\\
 \Lambda&=\bigoplus_{r=0}^K S V_r\subset E_S,\qquad
 \pi(x)=\sum_{r=0}^K\lambda_r(x)V_r,\\
 F&=\operatorname{span}_S
       \{y^{\mathbf b}:\mathbf b\notin\{\mathbf b^*(0),\ldots,
                                                   \mathbf b^*(K)\}\}.
 \end{split}
 \tag{TCL.8}\label{eq:tcl-splitting}
\]
Distinct \(V_r\) have distinct degrees, and the pivot coefficient
of \(V_r\) is the unit \(1/\mathbf b^*(r)!\). Therefore
\(\lambda_r(V_t)=\delta_{rt}\), proving both the direct sum and
\(\pi^2=\pi\). Subtracting \(\pi(x)\) deletes exactly all pivot
coordinates, so \(x-\pi(x)\in F\). Conversely all pivot
coordinates vanish on \(F\); hence \(\pi(F)=0\). We have proved
the explicit splitting
\[
 E_S=\Lambda\oplus F,\qquad \ker\pi=F.
 \tag{TCL.9}\label{eq:tcl-direct}
\]
In particular each \(V_r\) is primitive in the sense that it is
part of the displayed basis of a direct summand, with no division
by \(r!\) in the lattice construction.

Write \(L=f(\mathcal C_S)\). Equations (TCL.7)--(TCL.9) imply
\[
 \begin{split}
 L&=\bigoplus_{r=0}^K r!S V_r\subset\Lambda,\qquad
 \Lambda/L\simeq\bigoplus_{r=0}^K S/(r!),\\
 W_r&=U V_r,\quad \Lambda_U=U\Lambda,
 \quad L_U=f_U(\mathcal C_S)=UL,\quad F_U=UF,\\
 E_S&=\Lambda_U\oplus F_U,\qquad
 \pi_U=U\pi U^{-1},\qquad
 L_U=\bigoplus_{r=0}^K r!S W_r.
 \end{split}
 \tag{TCL.10}\label{eq:tcl-unit-lattices}
\]
Since \(S\) is a characteristic-zero domain, each \(r!\) is
nonzero and multiplication by it is injective. Thus \(f\) and
\(f_U\) are injective. The quotient isomorphism sends
\(\sum_r c_r V_r\) to \((c_r\bmod r!)_r\); its kernel and
surjectivity follow coordinate by coordinate. Applying the
automorphism \(U\) proves every transported identity in the
display. There is no claim that \(U\Lambda=\Lambda\); it is the
pair of lattices \(L\subset\Lambda\) that is transported.

Let \(Q=\operatorname{Frac}(S)\). The exact saturation statement is
\[
 (L_U\otimes_S Q)\cap E_S=\Lambda_U
       \quad\hbox{inside }E_S\otimes_S Q.
 \tag{TCL.11}\label{eq:tcl-saturation}
\]
Indeed the nonzero \(r!\) are invertible in \(Q\), so
\(L_U\otimes Q=\Lambda_U\otimes Q\). In the splitting
\(E_S=\Lambda_U\oplus F_U\), an integral element belonging
to this rational subspace has zero \(F_U\)-component, because
the free module \(F_U\) injects into \(F_U\otimes Q\).
This proves the stated intersection, including both inclusions.

Both lattices are stable under \(J\). On their indicated bases,
the full action and its covariant inclusion are
\[
 \begin{split}
 J W_r&=(r+1)W_{r+1},\qquad
 J(r!W_r)=(r+1)!W_{r+1},\\
 \mathcal C_S\xrightarrow{\ f_U\ }\Lambda_U,
 &\qquad T^r\longmapsto r!W_r,\qquad
 f_U\circ M_T=M_J\circ f_U.
 \end{split}
 \tag{TCL.12}\label{eq:tcl-action-matrix}
\]
At \(r=K\) both actions are zero. Commutation of \(U\) and
\(J\) and (TCL.7) prove all entries. This displays the actual
source multiplication and the containing-lattice multiplication
together with their exact intertwining map.

For clarity, the same map in all original monomial coordinates is
\[
 \begin{split}
 P_{ab}&=\begin{cases}\binom ba\rho^{b-a}&a\leq b,\\0&a>b,
                 \end{cases}\quad(0\leq a,b<m),\qquad P_k=P^{\otimes k},\\
 B_{rj}&=\begin{cases}\binom jr(k\rho)^{j-r}&r\leq j,\\0&r>j,
                 \end{cases}\quad(0\leq r,j\leq K),\\
 \mathsf V_{\mathbf b r}&=
       \begin{cases}1/\mathbf b!&|\mathbf b|=r,\\0&\text{otherwise},
       \end{cases}\qquad \Delta=\operatorname{diag}(0!,1!,\ldots,K!),\\
 [f_U^{\rm orig}]_{(s^{\mathbf b}),(Z^j)}
     &=P_k^{-1}\,[U]_{(y^{\mathbf b})}\,\mathsf V\,\Delta\,B.
 \end{split}
 \tag{TCL.13}\label{eq:tcl-original-matrix}
\]
The binomial theorem gives the two translation matrices \(P_k\)
and \(B\). Their inverses replace \(\rho\) and \(k\rho\) by
their negatives, respectively, as verified by the inverse substitutions
above; their diagonal entries are one. The \(j\)-th column on the
right is consequently the unchanged polynomial
\(U\sum_{r=0}^j\binom jr(k\rho)^{j-r}r!V_r=U(k\rho+J)^j\)
written in the original \(s^{\mathbf b}\)-basis. This proves the
matrix formula without replacing the original source map or its unit.

\subsection{Exact integer depths, quotient coordinates, and reduction}

For each integer \(r\geq0\) define
\[
 d_r=v_p(r!)=\sum_{a\geq1}\left\lfloor\frac{r}{p^a}\right\rfloor,
 \qquad r!=p^{d_r}\varepsilon_r,\qquad
 \varepsilon_r\in\mathbb Z_{(p)}^\times\subset S^\times.
 \tag{TCL.14}\label{eq:tcl-depth}
\]
Here \(0!=1\), and the sum is finite because its terms vanish
when \(p^a>r\). Each integer \(j\) between one and \(r\)
contributes one factor of \(p\) for each positive \(a\) with
\(p^a\mid j\). Interchanging these two finite counts yields
\(\sum_a\#\{1\leq j\leq r:p^a\mid j\}
=\sum_a\lfloor r/p^a\rfloor\), proving the valuation formula.
The remaining integer factor is prime to \(p\), hence is a unit
in \(S\). Therefore \((r!)=p^{d_r}S\) exactly. The depth cannot
increase in \(S\): an equality \(p^d=p^{d+1}s\) would permit
cancellation in the domain and give \(1=ps\), contradicting
\(\widetilde\alpha(1)=1\) and \(\widetilde\alpha(p)=0\).
This argument proves strictness of every consecutive pair of
principal \(p\)-power ideals, without assigning a valuation to
arbitrary elements of \(S\).

Set \(C=\Lambda_U/L_U\). The following are actual quotient maps:
\[
 \begin{split}
 q_r:E_S&\longrightarrow S/(p^{d_r}),\qquad
 q_r(x)=\lambda_r(U^{-1}x)\bmod p^{d_r},\\
 C&\xrightarrow{\ \sim\ }\bigoplus_{r=0}^K S/(p^{d_r}),
       \qquad [\sum_r c_rW_r]\longmapsto
                       (c_r\bmod p^{d_r})_r,\\
 E_S/L_U&\xrightarrow{\ \sim\ }
                 F_U\oplus\bigoplus_{r=0}^K S/(p^{d_r}).
 \end{split}
 \tag{TCL.15}\label{eq:tcl-quotient}
\]
The convention is \(S/(p^0)=0\). Each \(q_r\) is linear,
vanishes on \(F_U\) and on \(L_U\), and takes \(W_t\) to
\(\delta_{rt}\) in its indicated target. Equations
(TCL.9)--(TCL.10) prove the two isomorphisms and all these claims.
For the coordinate class \(\nu_r=[W_r]\in C\), its exact
annihilator in \(S\) is \(p^{d_r}S\), including the unit ideal
when \(\nu_r=0\) and \(d_r=0\).

Let \(B\) denote either \(S/pS\) with its quotient map, or
the specified target \(k_0\) with the map
\(\widetilde\alpha\). In both cases \(B\ne0\) and has
characteristic \(p\). Append a bar to each base-changed module,
element and map. Equations (TCL.8)--(TCL.10) base change directly,
because their projections and inverse changes of basis have their
coefficients in \(S\). Thus \(\overline W_0,\ldots,
\overline W_K\) remain a basis of the direct summand
\(\overline\Lambda_U\subset\overline E\). The original map becomes
\[
 \begin{gathered}
 \overline f_U(\overline T^r)=
 \begin{cases}
  \overline{r!}\,\overline W_r&0\leq r<p,\\
  0&p\leq r\leq K,
 \end{cases}\\
 \begin{aligned}
 \ker\overline f_U
    &=\bigoplus_{r=p}^K B\overline T^r=(\overline T^p),\\
 \operatorname{im}\overline f_U
    &=\bigoplus_{r=0}^{\min(K,p-1)}B\overline W_r.
 \end{aligned}
 \end{gathered}
 \tag{TCL.16}\label{eq:tcl-reduced-map}
\]
An empty direct sum is zero, and the ideal \((\overline T^p)\)
is zero when \(p>K\). For \(r<p\), every factor of \(r!\)
is a nonzero element of the prime field and therefore a unit in
\(B\). For \(r\geq p\) the factorial has a factor \(p\)
and maps to zero. Applying these facts to the independent
\(\overline W_r\) proves exactly the displayed kernel and image.

The nilpotency indices of the indicated actions are
\[
 \begin{array}{c|c}
 \text{module and action}&\text{nilpotency index}\\ \hline
 \overline{\mathcal C}_S=B[T]/(T^{K+1}),\ M_T&K+1\\
 \operatorname{im}\overline f_U,\ M_{\overline J}
           &\min(p,K+1)\\
 \overline E,\ M_{\overline J}&\min(p,K+1)\\
 \overline\Lambda_U,\ M_{\overline J}&\min(p,K+1).
 \end{array}
 \tag{TCL.17}\label{eq:tcl-indices}
\]
The nilpotency index is the least positive integer whose power is
zero; thus it is one for the zero action on a nonzero module.
In the first row, \(T^K\ne0\) in its free monomial basis and
\(T^{K+1}=0\). On \(\overline E\), the binomial coefficients
\(\binom pi\) for \(0<i<p\) are divisible by \(p\): their
factorial expression has one factor of \(p\) in its numerator
and none in its denominator. The commuting-variable binomial
formula therefore gives
\(\overline J^p=\sum_i\overline y_i^p=0\), since \(p>m\).
Also \(\overline J^{K+1}=0\). If
\(d=\min(p-1,K)\), then
\(\overline J^d=\overline{d!}\,\overline V_d\ne0\), by its
unit pivot coordinate. This proves the third row. Its application
to \(\overline W_0=\overline U\) gives
\(\overline J^d\overline W_0
=\overline{d!}\,\overline W_d\ne0\), proving the second and
fourth rows together with their inherited upper bounds.
In particular the abstract reduced cyclic source has not acquired
the smaller index: its exact map to the target has lost the
directions in (TCL.16).

The reduced containing lattice has the more precise chain decomposition
\[
 \overline\Lambda_U=
 \bigoplus_{\substack{j\geq0\\jp\leq K}}
       \bigoplus_{t=0}^{\min(p-1,K-jp)} B e_{j,t},\qquad
 e_{j,t}=\left(\prod_{a=1}^t\overline{jp+a}\right)
                          \overline W_{jp+t}.
 \tag{TCL.18}\label{eq:tcl-divided-chains}
\]
The empty product is one. All displayed products are units because
\(1\leq a\leq p-1\). These are therefore changes of basis on
disjoint intervals of the basis \(\overline W_r\). Formula
(TCL.12) gives \(\overline J e_{j,t}=e_{j,t+1}\) except at the
last position of a displayed interval. There it is zero, either
because the next coefficient is \(\overline{(j+1)p}=0\), or
because the next index is \(K+1\). This proves each full chain,
including every chain that is absent from the image in (TCL.16).

The lost directions are also the exact connecting image in the
base-change sequence
\[
 \begin{gathered}
 0\longrightarrow\mathcal C_S\xrightarrow{\ f_U\ }
             \Lambda_U\longrightarrow C\longrightarrow0,\\
 0\longrightarrow\operatorname{Tor}_1^S(C,B)
    \xrightarrow{\ \partial\ }\overline{\mathcal C}_S
    \xrightarrow{\ \overline f_U\ }\overline\Lambda_U
    \longrightarrow C\otimes_S B\longrightarrow0,\\
 \operatorname{Tor}_1^S(C,B)=\bigoplus_{r=p}^K B\epsilon_r,
       \qquad \partial(\epsilon_r)=\overline T^r,
       \qquad C\otimes_S B\simeq
                    \bigoplus_{r=p}^K B[\overline W_r].
 \end{gathered}
 \tag{TCL.19}\label{eq:tcl-tor}
\]
Here \(\operatorname{Tor}_1\) can be computed directly, without
a flatness assertion about \(B\). The first line is a free
resolution of \(C\) of length one, whose differential in the
displayed bases is \(\Delta\). Tensoring this two-term complex
with \(B\) gives differential \(\overline\Delta\).
By definition its degree-one homology is
\(\operatorname{Tor}_1^S(C,B)\). The factorial calculation in
(TCL.16) gives its kernel with the displayed basis and gives its
cokernel with the displayed basis. The degree-one differential
was chosen as \(+f_U\), so the map of the kernel into the
degree-one free module is precisely
\(\epsilon_r\mapsto\overline T^r\), with no additional sign.
This proves the entire sequence and its maps explicitly.

For the top original vector the formula, including its full scalar, is
\[
 f_U(T^K)=\frac{K!}{((m-1)!)^k}\,a_0^k
                    \prod_{i=1}^k y_i^{m-1}.
 \tag{TCL.20}\label{eq:tcl-terminal}
\]
Only \(\mathbf b=(m-1,\ldots,m-1)\) has total degree \(K\),
so (TCL.7) gives its multinomial coefficient. Every nonconstant
term of \(U\) annihilates that monomial; its constant coefficient
is exactly \(a_0^k\), proving the formula. Its scalar is nonzero
in the actual characteristic-zero domain, and its exact factorial
depth is \(d_K\), because \((m-1)!\) and \(a_0\) are units.
The displayed vector maps to zero under the specified characteristic
\(p\) coefficient map exactly when \(K\geq p\). This statement
concerns the already specified local single-primary source. It does
not extend that specialization to a larger coefficient ring containing
additional inverse CRT denominators. A ring in which \(p\) is a
unit cannot map unitally to characteristic \(p\), since applying
such a map to \(p p^{-1}=1\) would give \(0=1\).

\subsection{The exact original chain-support map for these depths}

Fix an integer \(L\geq\max(1,d_K)\) and put
\[
 \begin{gathered}
 R_{p,L}=\mathbb Z/p^L\mathbb Z,\qquad S_L=S/p^L S,\\
 \iota_L:R_{p,L}\longrightarrow S_L,
       \quad [z]\longmapsto[z],\\
 I_a=p^aR_{p,L},\quad H_a=p^aS_L\quad(0\leq a\leq L).
 \end{gathered}
 \tag{TCL.21}\label{eq:tcl-chain-rings}
\]
The scalar map is well-defined and unital. It is injective.
Indeed, if an integer \(z\) not divisible by \(p^L\) has image
zero, write \(z=p^a u\) with \(0\leq a<L\) and \(u\) prime
to \(p\). Then \(p^a u=p^L s\) in \(S\). Cancellation and
inversion of \(u\) imply \(1=p^{L-a}s u^{-1}\), contradicting
the specified residue map. Thus precisely \(p^L\mathbb Z\)
maps to zero, proving injectivity. The same argument proves
\(p^a\notin p^{a+1}S_L\) for \(a<L\): a purported equality
lifts to \(p^a=p^{a+1}s+p^L t\), and cancellation gives
\(1=p s+p^{L-a}t\), again impossible in the residue field.
Consequently all \(H_a\) are distinct.

For completeness, every ideal of \(R_{p,L}\) is some \(I_a\).
For a nonzero ideal, choose a nonzero element with least exponent
of \(p\) among its integer representatives. It has the form
\(p^a u\) with \(u\) a unit modulo \(p^L\), by the integer
Euclidean algorithm. Multiplying by \(u^{-1}\) puts \(p^a\)
in the ideal, and all its other elements have exponent at least
\(a\) by choice, proving that the ideal is \(I_a\). The zero
ideal is \(I_L\). This argument is not applied to \(S_L\):
there we retain exactly the principal subchain \(\{H_a\}\).

Extension and contraction on this actual subchain are
\[
 \iota_L(I_a)S_L=H_a,\qquad \iota_L^{-1}(H_a)=I_a,
 \qquad(0\leq a\leq L).
 \tag{TCL.22}\label{eq:tcl-extend-contract}
\]
Extension follows because the image of the principal generator is
\(p^a\). For contraction, an integer of exponent \(b<a\)
belonging to \(p^aS_L\) would give
\(p^b u=p^a s+p^L t\). Cancelling \(p^b\) contradicts the
nonzero residue of the unit \(u\). Thus only the integers divisible
by \(p^a\) contract, and they all do, proving the second formula.

The original ideal laws are retained exactly:
\[
 \begin{array}{lll}
 I_a+I_b=I_{\min(a,b)},& I_a\cap I_b=I_{\max(a,b)},&
              I_a I_b=I_{\min(L,a+b)},\\
 H_a+H_b=H_{\min(a,b)},& H_a\cap H_b=H_{\max(a,b)},&
              H_a H_b=H_{\min(L,a+b)}.
 \end{array}
 \tag{TCL.23}\label{eq:tcl-ideal-laws}
\]
For sum and intersection the proof is inclusion: when \(a\leq b\),
the ideal indexed by \(b\) is contained in the one indexed by
\(a\). Their sum is the larger and their intersection the smaller.
For product, every sum of products of multiples of \(p^a\) and
\(p^b\) is a multiple of \(p^{a+b}\), and the generator
\(p^{a+b}\) itself is such a product. It becomes zero for
\(a+b\geq L\). These arguments apply in both displayed rings
and prove all six laws, including the boundary indices.

Let \(C_L=\{0<1<\cdots<L\}\). The original lattice map and its
actual coefficient extension are
\[
 \theta_p:C_L\longrightarrow\operatorname{Id}(R_{p,L}),
       \quad j\longmapsto I_{L-j},\qquad
 \theta_S:C_L\longrightarrow\{H_a:0\leq a\leq L\},
       \quad j\longmapsto H_{L-j}.
 \tag{TCL.24}\label{eq:tcl-theta}
\]
They are bijections by the distinctness and classification just
proved. Each takes bottom to the zero ideal, top to the unit ideal,
maximum to ideal sum and minimum to ideal intersection, by
(TCL.23). Thus they are bounded-lattice isomorphisms onto their
stated codomains. Extension \(I_a\mapsto H_a\) commutes with
these two maps and preserves all three ideal operations of
(TCL.23). In the \(j\)-coordinates, ordinary ideal multiplication
is \(j\mathbin{\odot}t=\max(0,j+t-L)\), since its exponent is
\(\min(L,(L-j)+(L-t))\). It is not identified with lattice
meet: for \(L\geq2\), \(j=t=L-1\) gives respective values
\(L-2\) and \(L-1\).

For either coefficient ring \(B=R_{p,L}\) or \(B=S_L\), use
precisely the original split-support carrier
\[
 \begin{split}
 G_{C_L}(B)&=\{(0_B,j):0\leq j\leq L\}
                 \cup\{(b,L):b\in B\},\\
 (b,j)+(c,t)&=(b+c,\max(j,t)),\qquad
 (b,j)(c,t)=(bc,\min(j,t)),\\
 \tau_a&=(0_B,L-a)\ (0\leq a\leq L),\qquad
 e=\tau_0=(0_B,L),\qquad \tau_L=(0_B,0).
 \end{split}
 \tag{TCL.25}\label{eq:tcl-carrier}
\]
The union is inside \(B\times C_L\), so the two occurrences
of \((0_B,L)\) give one element. A nonzero amplitude occurs only
at top support. A nonzero sum has a nonzero summand, so has top
support; a nonzero product has both amplitudes nonzero, so has
top support. This proves closure. Associativity and commutativity
hold in each coordinate. Distributivity follows from ring
distributivity and \(\min(j,\max(t,v))
=\max(\min(j,t),\min(j,v))\), verified by ordering \(t,v\).
The zero is \(\tau_L\), and the unit is \((1_B,L)\).
These facts prove the stated semiring structure directly.

The exact support-preserving coefficient morphism is
\[
 \begin{split}
 G_{C_L}(\iota_L):G_{C_L}(R_{p,L})&\longrightarrow G_{C_L}(S_L),
           \quad (b,j)\longmapsto(\iota_L(b),j),\\
 \chi_B:G_{C_L}(B)&\longrightarrow C_L,
           \quad(b,j)\longmapsto j,\\
 \theta_S\chi_{S_L}\,G_{C_L}(\iota_L)
    &=\operatorname{ext}_{\iota_L}\,\theta_p\chi_{R_{p,L}}.
 \end{split}
 \tag{TCL.26}\label{eq:tcl-support-map}
\]
The first map is well-defined because nonzero target amplitude
has nonzero source amplitude and hence top support. Applying the
ring homomorphism to amplitudes and leaving the support coordinate
fixed proves preservation of both operations, zero and unit.
Its injectivity follows from injectivity of \(\iota_L\) and
the unchanged support coordinate. The support projection preserves
the two operations by their defining formulas. The last equality
follows on every entry \((b,j)\) from
\(I_{L-j}S_L=H_{L-j}\); it is an equality into the principal
ideal subchain with ideal sum and ideal intersection. In particular
each \(\tau_a\) maps to the distinct \(\tau_a\), then to
\(I_a\) or \(H_a\). The supported zero \(e\) remains distinct
from the absent zero \(\tau_L\).

Finally the connection of this support diagram to the actual cyclic
quotient is not an identification of a torsion module with a semiring.
The precise coordinate and annihilator maps are
\[
 \begin{split}
 C&\simeq\bigoplus_{r=0}^K S/(p^{d_r}),\qquad p^L C=0,\\
 \operatorname{Ann}_{S_L}(\nu_r)&=H_{d_r}
     =\theta_S(L-d_r)
     =\theta_S\chi_{S_L}(\tau_{d_r}),\\
 q_r(x)&=\lambda_r(U^{-1}x)\bmod p^{d_r}.
 \end{split}
 \tag{TCL.27}\label{eq:tcl-quotient-support}
\]
The first and last lines are the proved maps of (TCL.15), with
\(p^L\) annihilating every summand because \(d_r\leq L\).
An element of \(S_L\) kills the unit class of the \(r\)-th
summand exactly when a representative in \(S\) lies in
\(p^{d_r}S\); this proves the annihilator formula and all its
equalities using (TCL.24)--(TCL.26). Thus the depths of the
original factorial inclusion have explicit labels in the original
chain-support carrier, with their arithmetic ideal product as well
as their lattice sum and intersection still recorded.
The entry \(\tau_{d_r}\) labels the annihilator; it does not
replace the vector \(\nu_r\). In particular, when \(d_r=L\)
the vector \(\nu_r\) is the nonzero unit class in \(S_L\), whereas
its annihilator is the zero ideal labelled by \(\tau_L\).

\subsection{The retained boundary-character labels and the full-stage obstruction}

The action on the lost source directions is also exact. For
\(K\geq p\), the kernel in (TCL.16), equivalently the connecting
image of the Tor group in (TCL.19), has the ordered basis
\(\overline T^p,\ldots,\overline T^K\), with
\[
 M_T\overline T^r=\overline T^{r+1}\ (p\leq r<K),\qquad
 M_T\overline T^K=0,\qquad
 \operatorname{ind}\bigl(M_T|_{\ker\overline f_U}\bigr)=K-p+1.
 \tag{TCL.28}\label{eq:tcl-tor-index}
\]
These equalities follow in the free monomial source basis. In
particular the \((K-p)\)-th power takes \(\overline T^p\)
to the nonzero \(\overline T^K\), and the next power is zero
on every basis element. This index can exceed \(p\). It is an
action on the exact kernel in the source, not a subspace of the
ambient target. Thus it is compatible with
\(\overline J^p=0\) on \(\overline E\): the original map is
zero on every one of these kernel vectors.

Put \(N=m+1\). For a residue class \(c\in\mathbb Z/N\mathbb Z\)
define the original coefficient character masks in \(y\)-coordinates
and their full-unit transports by
\[
 \begin{split}
 Q_c(y^{\mathbf b})&=
       \mathbf 1_{\{|\mathbf b|+k\equiv c\ ({\rm mod}\ N)\}}
                       y^{\mathbf b},\qquad
 \widehat Q_c=UQ_cU^{-1}:E_S\longrightarrow E_S,\\
 D_c(T^r)&=\mathbf 1_{\{r+k\equiv c\ ({\rm mod}\ N)\}}T^r
                       \quad\hbox{on }\mathcal C_S.
 \end{split}
 \tag{TCL.29}\label{eq:tcl-character-mask}
\]
Each diagonal entry defining \(Q_c,D_c\) is zero or one in
\(S\), so these maps exist over this exact coefficient ring
without adjoining roots of unity. On every indicated basis vector
exactly one of the masks is one. Hence each family consists of
pairwise orthogonal algebraic idempotents summing to the identity;
conjugation by the actual \(U\) proves the same assertion for
\(\widehat Q_c\). The word orthogonal here means that distinct
idempotents compose to zero; no inner product or orthogonal
projection on a metric space is introduced.

The full original boundary-character comparison is
\[
 \begin{split}
 Q_c V_r&=\mathbf 1_{\{r+k\equiv c\}}V_r,\qquad
 \widehat Q_c W_r=\mathbf 1_{\{r+k\equiv c\}}W_r,\\
 \widehat Q_c f_U&=f_U D_c,\qquad
 \widehat Q_c J=J\widehat Q_{c-1},\qquad
 D_c M_T=M_TD_{c-1}.
 \end{split}
 \tag{TCL.30}\label{eq:tcl-character-diagram}
\]
Every summand of \(V_r\) has degree \(r\), proving the first
equality, and conjugation proves the second. Apply these formulas
to \(f_U(T^r)=r!W_r\) to prove the third equality on a source
basis. Multiplication by each \(y_i\) raises degree by one for
every surviving monomial and gives zero at the truncation boundary.
It follows on the complete monomial basis that
\(Q_cJ=JQ_{c-1}\). Conjugation and \(UJ=JU\) prove the fourth
identity; the identical degree argument on \(T^r\) proves the
last. These are covariant character shifts, not a claim that a
fixed character mask commutes with \(J\).

Over \(\mathbb C\), this is exactly the coefficient mask of
the boundary monodromy in TP.9--10: that period monodromy has
eigenvalue \(\exp(2\pi i(|\mathbf b|+k)/N)\) on its period
character indexed by \(\mathbf b\). In the unchanged original
\(s^{\mathbf b}\)-coordinates the coefficient mask is
\(P_k^{-1}Q_cP_k\), and the source mask is its conjugate by
\(P_k^{-1}[U]_{(y^{\mathbf b})}P_k\). The complex period matrix
is \(\Pi_k=e^{kc_0/u}W_kD_k(u)P_k\), with the original phase
constant here denoted \(c_0=-(-\rho)^N/N\), and \(D_k(u)\)
diagonal in these same \(\mathbf b\) indices. Consequently
\[
 \mathcal P_c\Pi_k=\Pi_k(P_k^{-1}Q_cP_k),\qquad
 \mathcal P_c\,\Pi_k[U]_{(s^{\mathbf b})}^{-1}
   =\Pi_k[U]_{(s^{\mathbf b})}^{-1}
                [\widehat Q_c]_{(s^{\mathbf b})},
 \quad \mathcal P_c=W_kQ_cW_k^{-1}.
 \tag{TCL.31}\label{eq:tcl-period-mask}
\]
Indeed \(Q_c\) commutes with the diagonal \(D_k(u)\), which
proves the first equality by multiplication, and substitution of
the conjugation formula for \(\widehat Q_c\) proves the second.
The common scalar \(e^{kc_0/u}\), the original \(P_k\) and the
entire \(U\) are present in these identities.

The masks preserve both \(L_U\) and \(\Lambda_U\), as follows
from their diagonal action on \(r!W_r\) and \(W_r\). They
therefore act on the quotient \(C\), and the exact free resolution
in (TCL.19) commutes with \(D_c\) on its source and
\(\widehat Q_c\) on its target. Its base change and its kernel
inherit precisely those same masks. Thus the simultaneously retained
depth and character of every quotient coordinate and lost direction are
\[
 \begin{split}
 \nu_r=[W_r]&:\quad
       \operatorname{Ann}_{S_L}(\nu_r)=H_{d_r},\quad
       \widehat Q_c\nu_r=\mathbf 1_{\{r+k\equiv c\}}\nu_r,\\
 \epsilon_r\in\operatorname{Tor}_1^S(C,B)\quad(r\geq p)
     &:\quad \overline D_c\partial\epsilon_r
        =\mathbf 1_{\{r+k\equiv c\}}\overline T^r.
 \end{split}
 \tag{TCL.32}\label{eq:tcl-two-labels}
\]
The first line combines the proved annihilator in (TCL.27) with
the induced quotient map. The second follows from the explicit
connecting map \(\partial\epsilon_r=\overline T^r\).
This proves both labels on the actual maps, including the zero
coordinates when \(d_r=0\). It does not assert an inertia-invariant
description of a different sheaf after reduction.

There is an exact obstruction to carrying this local specialization
through the full amplified quartet CRT stage. In the counterfactual
quartet of PAM.1--4 retain
\(\rho=\tfrac12+\delta+i\gamma\), with
\(0<\delta<\tfrac12\), \(\gamma>0\), and take \(k\geq p\).
Two actual ordered tensor centres in that construction are
\[
 \lambda_0=k\rho,\qquad
 \lambda_1=(k-p)\rho+p(1-\overline\rho),\qquad
 \lambda_0-\lambda_1=p(\rho+\overline\rho-1)=2p\delta\ne0.
 \tag{TCL.33}\label{eq:tcl-crt-centres}
\]
The first chooses \(\rho\) in every factor; the second chooses
\(1-\overline\rho\) in exactly \(p\) factors and \(\rho\)
in the other \(k-p\) factors. Both choices are permitted in the
retained quartet, and the difference follows by subtraction.
In the grid coordinates of PAM.2 they have indices respectively
\((a,b)=(k,k)\) and \((k-p,k)\).

Let \(A_{\rm CRT}\subset\mathbb C\) be the actual coefficient
ring carrying those quartet coordinates and every coefficient of
the retained terminal CRT idempotent polynomial representative
\(e(Z)=\sum_{r=0}^{d}e_rZ^r\) at \(\lambda_0\).
By its specified CRT residues, \(e(\lambda_0)=1\) and
\(e(\lambda_1)=0\). Define the following element of this same ring:
\[
 \mathscr D=\sum_{r=1}^{d}e_r
       \sum_{j=0}^{r-1}\lambda_0^{r-1-j}\lambda_1^j
       \in A_{\rm CRT},\qquad
 1=p(\rho+\overline\rho-1)\mathscr D,
 \qquad
 p^{-1}=(\rho+\overline\rho-1)\mathscr D\in A_{\rm CRT}.
 \tag{TCL.34}\label{eq:tcl-crt-obstruction}
\]
For every \(r\geq1\), multiplying the inner sum by
\(\lambda_0-\lambda_1\) cancels all intermediate monomials
and leaves \(\lambda_0^r-\lambda_1^r\). After multiplication
by \(e_r\) and summation the result is
\(e(\lambda_0)-e(\lambda_1)=1\). Substitution of (TCL.33)
proves the second equality and then the displayed inverse of
\(p\). This proof uses the original CRT polynomial itself and
does not depend on a selected formula for its coefficients or
on an assumed list of inverted pairwise differences.
Since \(p\) is a unit, no prime ideal of \(A_{\rm CRT}\)
lies over \((p)\), and there is no unital characteristic-\(p\)
specialization of that complete stage. A prime containing \(p\)
would contain \(1\), and any such specialization would send the
proved identity \(p p^{-1}=1\) to \(0=1\).

Thus for \(k\geq p\) the actual local prime used in (TCL.1)
does not extend to this full quartet stage with its retained
terminal CRT idempotent. This is a proved obstruction for that
specified stage, not a statement that the local and global
characteristic-zero objects are unrelated. Their local cyclic
coefficient map is still exactly (TCL.4), their unchanged monomial
translation is (TCL.13), and the characteristic-zero terminal
coefficient is (TCL.20). The local characteristic-\(p\) quotient
and its lost directions are the maps (TCL.16)--(TCL.19); a missing
extension of their coefficient map is not the image of a global
vector under a nonexistent specialization.

More generally let \(A_{{\rm CRT},k}\subset\mathbb C\) denote the
retained coefficient ring of the complete split corner at degree
\(k\). The argument (TCL.33)--(TCL.34), applied to every rational
prime \(q\leq k\) in place of \(p\), proves
\[
 q^{-1}\in A_{{\rm CRT},k}\quad(q\text{ prime},\ q\leq k).
 \tag{TCL.35}\label{eq:tcl-all-small-primes}
\]
Indeed the second centre uses \(q\) reflected factors, which are
available for every such \(q\), and the same retained idempotent
separates it from the terminal centre. Its divided difference
gives the inverse without changing the idempotent.
For a directed coefficient ring containing all these actual
split stages one may take the explicit subring
\[
 A_{{\rm CRT},\infty}
   =\bigcup_{n\geq2}
       \mathbb Z[A_{{\rm CRT},2},\ldots,A_{{\rm CRT},n}]
       \ \subset\mathbb C,\qquad
 \mathbb Q\subset A_{{\rm CRT},\infty}.
 \tag{TCL.36}\label{eq:tcl-split-limit}
\]
The finite generated-stage subrings in this union form an
increasing sequence, so the union is a subring containing every
stage with its actual coefficient maps. For any prime \(q\),
the stage \(k=q\) supplies \(q^{-1}\) by (TCL.35).
Factoring a nonzero integer into primes then puts its reciprocal
in the union, proving the asserted copy of \(\mathbb Q\).
Every residue field of this union therefore has characteristic
zero: a characteristic-\(q\) ring homomorphism would send the
unit identity \(q q^{-1}=1\) to \(0=1\).
Equivalently the scalar morphism
\(\operatorname{Spec}A_{{\rm CRT},\infty}
 \to\operatorname{Spec}\mathbb Z\)
has no point over any \((q)\). This follows directly by contraction
of primes and the proved invertibility of every \(q\).
These conclusions concern the full split corner diagram with
its retained CRT idempotents. They do not remove the finite
single-primary models or the local maps (TCL.1)--(TCL.32).

There are also exact surviving character-zero vectors in the
reduced local cyclic image. Take \(m\geq2\) and \(k=p\ell\) with
\(\ell\geq1\), so \(K=k(m-1)\geq p\). Its final nonzero cyclic
image basis vector and its unchanged characteristic-zero precursor
satisfy
\[
 \begin{split}
 \overline J\,\overline W_{p-1}&=0,\qquad
       \overline W_{p-1}\ne0,\qquad J W_{p-1}=pW_p\ne0,\\
 \overline{\widehat Q}_0\,\overline W_{p-1}=\overline W_{p-1}
       &\quad\Longleftrightarrow\quad
       p(\ell+1)\equiv1\pmod{N},\qquad N=m+1 .
 \end{split}
 \tag{TCL.37}\label{eq:tcl-surviving-character}
\]
The first line follows from (TCL.12), the unit pivot after
reduction, and the nonzero scalar \(p\) on the free
characteristic-zero containing lattice. The vector belongs to
the original reduced cyclic image because
\(\overline f_U(\overline T^{p-1})
 =\overline{(p-1)!}\,\overline W_{p-1}\), whose scalar is a
unit. Formula (TCL.30) gives the second line as
\(p-1+k\equiv0\pmod N\), and substitution of \(k=p\ell\)
gives the stated congruence. In the retained coefficient datum
\(p>N\), so \(p\) is invertible modulo \(N\). Choosing the
single residue class
\(\ell\equiv p^{-1}-1\pmod N\) and then its infinitely many
positive representatives proves that the displayed character-zero
case occurs for infinitely many tensor orders \(k=p\ell\).
These vectors are killed by \(J\) only after base change, as
their exact characteristic-zero image \(pW_p\) shows. This is
a calculation in the local coefficient module with its retained
character masks, not an identification with a finite-field
sheaf vector or a recovery of the original terminal jet.

Every tensor coefficient, the original monomial translations, the
full unit and inverse, and the specified coefficient map occur in
the displayed morphisms. The characteristic-\(p\) kernel in
(TCL.16) is a statement about these coefficient modules. It neither
sets an integer scalar to zero in \(\overline{\mathbb Q}_\ell\)
nor identifies this reduction with the inertia or Frobenius action
on the finite-field sheaf. No verdict about the location of an
actual zero of \(\xi\) follows from this lattice calculation alone.
